% =============================================================================
% APPENDIX PR: CORE-PRINCIPLES: Fundamental Axioms of Evidence-Based Frameworks
% =============================================================================
% Module: BCM2_00_META
% Version: 1.0 (January 2026)
% Status: Draft
%
% This appendix defines the 7 fundamental principles (axioms) that constitute
% an Evidence-Based Framework (EBF) for economic and social behavior.
%
% =============================================================================

% =============================================================================
% CROSS-REFERENCE STRUCTURE
% =============================================================================
%
% CHAPTER LINKAGE (Required):
% - Primary Chapter: Chapter 1 (Introduction to EBF)
% - Secondary Chapters: All chapters (principles apply throughout)
%
% DEPENDENCIES (what this appendix REQUIRES):
% - Appendix G: Glossary (symbol definitions)
%
% DEPENDENTS (what REQUIRES this appendix):
% - All other EBF appendices (inherit these principles)
% - Appendix UNMAPPED_HI: CORE-INTELLIGENCE (Menschlichkeit × Mechanik)
% - Appendix UNMAPPED_BX: CORE-BEATRIX (three-tier hierarchy)
% - Appendix UNMAPPED_FM: METHOD-FRAMEWORK (scientific reusability)
%
% =============================================================================

\begin{tcolorbox}[colback=purple!5!white, colframe=purple!75!black,
    title=Chapter Linkage]

\textbf{Primary Chapter:} Chapter 1: Introduction to EBF
\begin{itemize}[nosep]
\item Section 1.1: What is EBF? $\leftrightarrow$ This Appendix Section \ref{sec:pr-overview}
\item Section 1.2: Framework vs. Model $\leftrightarrow$ This Appendix Section \ref{sec:pr-principles}
\end{itemize}

\textbf{Secondary Chapters:}
\begin{itemize}[nosep]
\item Chapter 5: Utility (WHAT) --- P2, P3
\item Chapter 9: Context (WHEN) --- P2
\item Chapter 10: Complementarity (HOW) --- P1
\end{itemize}

\textbf{Reading Guide:}
\begin{itemize}[nosep]
\item For overview: This appendix first (foundational)
\item For detailed theory: Follow to specific CORE appendices (AAA, B, C, V)
\item For implementation: Appendix UNMAPPED_BX (BEATRIX architecture)
\end{itemize}

\end{tcolorbox}

\chapter{CORE-PRINCIPLES: Fundamental Axioms of Evidence-Based Frameworks}
\label{app:core-principles}


\begin{tcolorbox}[colback=blue!5!white, colframe=blue!75!black,
    title=Quick Reference: Appendix UNMAPPED_PR]

\textbf{Appendix:} PR (CORE-PRINCIPLES)

\textbf{Core Question:} What are the fundamental principles that define an Evidence-Based Framework for economic and social behavior?

\textbf{Key Formula:}
\begin{equation}
\text{EBF} = \underbrace{(P_1 + P_2 + P_3)}_{\text{Behavioral Theory}} \times \underbrace{(P_4 + P_5)}_{\text{Scientific Method}} \times \underbrace{(P_6 + P_7)}_{\text{Scalability}}
\label{eq:pr-ebf}
\end{equation}

\textbf{Key Concepts:}
\begin{itemize}[nosep]
\item \textbf{7 Principles (P1--P7)}: The axiomatic foundation of EBF
\item \textbf{Complementarity ($\gamma$)}: Dimensions interact, not substitute
\item \textbf{Context ($\Psi$)}: Situation shapes behavior
\item \textbf{Menschlichkeit}: What machines cannot replicate
\end{itemize}

\textbf{Cross-References:} Appendix UNMAPPED_AAA (WHO), B (HOW), C (WHAT), V (WHEN), HI (INTELLIGENCE), BX (BEATRIX), G (Glossary)
\end{tcolorbox}


\begin{abstract}
This appendix formalizes the \textbf{seven fundamental principles} that constitute an Evidence-Based Framework (EBF) for economic and social behavior. These principles---Complementarity (P1), Context (P2), Systematik (P3), Evidence (P4), Falsifiability (P5), Human-Machine Complementarity (P6), and Cumulation (P7)---define what distinguishes EBF from ad-hoc behavioral analysis. Together, they form a coherent axiomatic system that guides all EBF applications.
\end{abstract}


\begin{tcolorbox}[
    colback=orange!5!white,
    colframe=orange!75!black,
    title={\textbf{Appendix Scope: Ziel / In-Scope / Out-of-Scope / Constraints}},
    fonttitle=\bfseries
]

\textbf{Ziel (Objective):}
\begin{itemize}[nosep]
    \item Define the 7 fundamental axioms that constitute an Evidence-Based Framework for behavioral analysis
\end{itemize}

\textbf{In-Scope:}
\begin{itemize}[nosep]
    \item Formal statement of 7 principles (P1--P7) as axioms
    \item Mathematical formalization of each principle
    \item Interdependencies between principles
    \item Cross-references to implementing appendices
\end{itemize}

\textbf{Out-of-Scope (delegiert an andere Appendices):}
\begin{itemize}[nosep]
    \item Detailed 10C framework $\rightarrow$ Appendix UNMAPPED_AAA-HI (CORE appendices)
    \item Technical LLM architecture $\rightarrow$ Appendix UNMAPPED_BX (CORE-BEATRIX)
    \item Mechanics Thesis details $\rightarrow$ Appendix UNMAPPED_HI (CORE-INTELLIGENCE)
    \item Parameter estimation $\rightarrow$ Appendix UNMAPPED_BBB (CORE-WHERE)
\end{itemize}

\textbf{Constraints (Anwendungsgrenzen):}
\begin{itemize}[nosep]
    \item Principles are framework-level, not model-level
    \item Principles are necessary but not sufficient for complete EBF
    \item Implementation requires all 7 principles together
\end{itemize}

\textbf{Lieferobjekte (Deliverables):}
\begin{itemize}[nosep]
    \item 7 formal axioms (UNMAPPED_PR-1 to UNMAPPED_PR-7)
    \item Principle interdependency matrix
    \item EBF completeness theorem
    \item Worked example applying all 7 principles
\end{itemize}

\end{tcolorbox}


%%%%%%%%%%%%%%%%%%%%%%%%%%%%%%%%%%%%%%%%%%%%%%%%%%%%%%%%%%%%%%%%%%%%%%%%%%%%%%%
\section{The Fundamental Question}
\label{sec:pr-fundamental}
%%%%%%%%%%%%%%%%%%%%%%%%%%%%%%%%%%%%%%%%%%%%%%%%%%%%%%%%%%%%%%%%%%%%%%%%%%%%%%%

\subsection{Why Principles Matter}

Many approaches to behavioral analysis exist: traditional economics, behavioral economics, psychology, sociology, marketing science. Each offers insights, but none provides a \textit{systematic, evidence-based, scalable} framework for analyzing economic and social behavior.

\begin{tcolorbox}[colback=red!5!white, colframe=red!75!black,
    title=The Fundamental Question]
\textbf{What are the minimal axioms required for a rigorous, evidence-based, scalable framework for analyzing human behavior?}
\end{tcolorbox}

\subsection{What Distinguishes EBF}

An Evidence-Based Framework differs from:

\begin{table}[htbp]
\centering
\caption{EBF vs. Alternative Approaches}
\label{tab:pr-comparison}
\renewcommand{\arraystretch}{1.4}
\begin{tabular}{@{}p{3cm}p{4cm}p{5cm}@{}}
\toprule
\textbf{Approach} & \textbf{Limitation} & \textbf{EBF Addresses Via} \\
\midrule
Traditional Econ & Single utility, no context & P1 (Complementarity), P2 (Context) \\
Ad-hoc Behavioral & No systematic structure & P3 (10C Systematik) \\
Pure LLM Chat & No evidence, no validation & P4 (Evidence), P5 (Falsifiability) \\
Manual Consulting & Doesn't scale & P6 (Human-Machine), P7 (Cumulation) \\
\bottomrule
\end{tabular}
\end{table}


%%%%%%%%%%%%%%%%%%%%%%%%%%%%%%%%%%%%%%%%%%%%%%%%%%%%%%%%%%%%%%%%%%%%%%%%%%%%%%%
\section{Overview: The 7 Principles}
\label{sec:pr-overview}
%%%%%%%%%%%%%%%%%%%%%%%%%%%%%%%%%%%%%%%%%%%%%%%%%%%%%%%%%%%%%%%%%%%%%%%%%%%%%%%

\begin{tcolorbox}[colback=green!10!white, colframe=green!75!black,
    title=\textbf{The 7 Fundamental Principles of EBF}]

\begin{enumerate}
    \item[\textbf{P1}] \textbf{COMPLEMENTARITY} --- People have multiple goals that interact
    \item[\textbf{P2}] \textbf{CONTEXT} --- Situation shapes behavior more than personality
    \item[\textbf{P3}] \textbf{SYSTEMATIK} --- Use structured 10C analysis, not ad-hoc
    \item[\textbf{P4}] \textbf{EVIDENCE} --- Base claims on scientific literature, not opinion
    \item[\textbf{P5}] \textbf{FALSIFIABILITY} --- Make testable predictions with confidence intervals
    \item[\textbf{P6}] \textbf{HUMAN-MACHINE} --- Combine Menschlichkeit with mechanical amplification
    \item[\textbf{P7}] \textbf{CUMULATION} --- Build knowledge that accumulates over time
\end{enumerate}

\vspace{0.5em}
\textbf{The EBF Equation:}
\begin{equation}
\text{EBF} = \underbrace{(P_1 + P_2 + P_3)}_{\text{Behavioral Theory}} \times \underbrace{(P_4 + P_5)}_{\text{Scientific Method}} \times \underbrace{(P_6 + P_7)}_{\text{Scalability}}
\end{equation}

\textbf{Key Insight:} All 7 principles are necessary. Removing any one collapses the framework.

\end{tcolorbox}


%%%%%%%%%%%%%%%%%%%%%%%%%%%%%%%%%%%%%%%%%%%%%%%%%%%%%%%%%%%%%%%%%%%%%%%%%%%%%%%
\section{The Seven Principles: Formal Axioms}
\label{sec:pr-principles}
%%%%%%%%%%%%%%%%%%%%%%%%%%%%%%%%%%%%%%%%%%%%%%%%%%%%%%%%%%%%%%%%%%%%%%%%%%%%%%%

%-----------------------------------------------------------------------
\subsection{Principle 1: Complementarity ($\gamma$)}
%-----------------------------------------------------------------------

\begin{tcolorbox}[colback=gray!5!white, colframe=gray!75!black,
    title=Axiom UNMAPPED_PR-1: Complementarity]

\textbf{Statement:} Humans have multiple utility dimensions that \textit{interact}, not merely sum:

\begin{equation}
U = \sum_{d \in D} u_d + \sum_{d \neq d'} \gamma_{dd'} \cdot u_d \cdot u_{d'}
\label{eq:pr-complementarity}
\end{equation}

where:
\begin{itemize}[nosep]
    \item $D = \{\text{F, E, P, S, D, E}\}$ (Financial, Ego, Psychological, Social, Domain, Existential)
    \item $u_d$ = utility from dimension $d$
    \item $\gamma_{dd'}$ = complementarity coefficient between dimensions $d$ and $d'$
\end{itemize}

\textbf{Interpretation:}
\begin{itemize}[nosep]
    \item $\gamma > 0$: Dimensions are \textit{complements} (more X makes Y more valuable)
    \item $\gamma < 0$: Dimensions are \textit{substitutes} (more X makes Y less valuable)
    \item $\gamma = 0$: Dimensions are \textit{independent}
\end{itemize}

\textbf{Implication:} Interventions targeting one dimension can affect others. The $\gamma$-matrix determines spillover effects.

\textbf{Example:} Blood donation + monetary payment: $\gamma_{\text{FS}} < 0$ (payment crowds out social motivation).

\textbf{Cross-Reference:} Appendix UNMAPPED_B (CORE-HOW) for full $\gamma$-matrix specification.

\end{tcolorbox}


%-----------------------------------------------------------------------
\subsection{Principle 2: Context ($\Psi$)}
%-----------------------------------------------------------------------

\begin{tcolorbox}[colback=gray!5!white, colframe=gray!75!black,
    title=Axiom UNMAPPED_PR-2: Context Primacy]

\textbf{Statement:} Behavior is determined by the interaction of person and situation:

\begin{equation}
\text{Behavior} = f(\text{Person}, \Psi)
\end{equation}

where $\Psi = (\Psi_I, \Psi_S, \Psi_C, \Psi_K, \Psi_E, \Psi_T, \Psi_M, \Psi_F)$ represents 8 context dimensions:

\begin{table}[H]
\centering
\renewcommand{\arraystretch}{1.3}
\begin{tabular}{@{}cll@{}}
\toprule
\textbf{Symbol} & \textbf{Dimension} & \textbf{Description} \\
\midrule
$\Psi_I$ & Institutional & Rules, laws, defaults \\
$\Psi_S$ & Social & Who is present, social norms \\
$\Psi_C$ & Cognitive & Mental state (tired, stressed, attentive) \\
$\Psi_K$ & Cultural & Values, traditions, beliefs \\
$\Psi_E$ & Economic & Resources (budget, time, energy) \\
$\Psi_T$ & Temporal & When (time of day, life stage, urgency) \\
$\Psi_M$ & Material & Technology, infrastructure, tools \\
$\Psi_F$ & Physical & Location (home, office, public) \\
\bottomrule
\end{tabular}
\end{table}

\textbf{Interpretation:} The same person makes different decisions in different contexts. Context often explains more variance than personality.

\textbf{Implication:} Before analyzing behavior, always characterize $\Psi$. Without context, predictions are unreliable.

\textbf{Cross-Reference:} Appendix UNMAPPED_V (CORE-WHEN) for full context specification.

\end{tcolorbox}


%-----------------------------------------------------------------------
\subsection{Principle 3: 10C Systematik}
%-----------------------------------------------------------------------

\begin{tcolorbox}[colback=gray!5!white, colframe=gray!75!black,
    title=Axiom UNMAPPED_PR-3: Systematic 10C Analysis]

\textbf{Statement:} Every behavioral question must be analyzed through 10 CORE dimensions:

\begin{equation}
\text{Analysis} = \bigcap_{i=1}^{10} \text{CORE}_i
\end{equation}

\begin{table}[H]
\centering
\renewcommand{\arraystretch}{1.2}
\begin{tabular}{@{}clll@{}}
\toprule
\textbf{\#} & \textbf{CORE} & \textbf{Question} & \textbf{Output} \\
\midrule
1 & WHO & Who has utility? & Levels $L$ \\
2 & WHAT & What drives utility? & Dimensions $d \in D$ \\
3 & HOW & How do dimensions interact? & $\gamma$-matrix \\
4 & WHEN & When does context matter? & $\Psi$-vector \\
5 & WHERE & Where do parameters come from? & $\Theta$ (estimates) \\
6 & AWARE & How aware is the person? & $A(\cdot)$ \\
7 & READY & How ready to act? & $W_{AX}, \tau, \theta$ \\
8 & STAGE & Where in the journey? & $\phi$, $S(t)$ \\
9 & HIERARCHY & How do decisions stratify? & $L_0$--$L_3$ \\
10 & EIT & How do interventions emerge? & $\vec{I} \in [0,1]^9$ \\
\bottomrule
\end{tabular}
\end{table}

\textbf{Interpretation:} The 10C framework ensures no critical dimension is overlooked. Incomplete analysis leads to flawed predictions.

\textbf{Implication:} Every EBF analysis must address all 10 COREs, even if some are trivial for a given case.

\textbf{Cross-Reference:} Appendices AAA, B, C, V, BBB, AU, AV, AW, HI, IE (one for each CORE).

\end{tcolorbox}


%-----------------------------------------------------------------------
\subsection{Principle 4: Evidence-Based}
%-----------------------------------------------------------------------

\begin{tcolorbox}[colback=gray!5!white, colframe=gray!75!black,
    title=Axiom UNMAPPED_PR-4: Evidence Primacy]

\textbf{Statement:} Every claim must be supported by scientific evidence:

\begin{equation}
\forall \text{ claim } C: \exists \text{ evidence } E \text{ such that } P(C | E) > P(C)
\end{equation}

\textbf{Evidence Hierarchy:}
\begin{enumerate}
    \item \textbf{Tier 1 (Gold):} RCTs, Top-5 journals, replicated findings
    \item \textbf{Tier 2 (Silver):} Peer-reviewed, solid methodology (JF, MS, JEBO, NBER)
    \item \textbf{Tier 3 (Bronze):} Working papers, theoretical models
\end{enumerate}

\textbf{EBF Evidence Base:}
\begin{itemize}[nosep]
    \item 2,328 curated papers in \texttt{bcm\_master.bib}
    \item 153 formalized theories in \texttt{theory-catalog.yaml}
    \item 852 documented cases in \texttt{case-registry.yaml}
    \item 400+ empirically estimated parameters
\end{itemize}

\textbf{Interpretation:} Opinions, intuitions, and ``common sense'' are not evidence. Claims without evidence must be marked as hypotheses.

\textbf{Implication:} Before making a claim, search internal sources (bibliography, cases, theories) then external sources.

\textbf{Cross-Reference:} Evidence Integration Pipeline (EIP) in \texttt{docs/workflows/evidence-integration-pipeline.md}.

\end{tcolorbox}


%-----------------------------------------------------------------------
\subsection{Principle 5: Falsifiability}
%-----------------------------------------------------------------------

\begin{tcolorbox}[colback=gray!5!white, colframe=gray!75!black,
    title=Axiom UNMAPPED_PR-5: Falsifiable Predictions]

\textbf{Statement:} Every EBF prediction must be falsifiable with quantitative bounds:

\begin{equation}
\text{Prediction} = \hat{Y} \pm \sigma \quad \text{with} \quad P(\hat{Y} - k\sigma \leq Y \leq \hat{Y} + k\sigma) = 1 - \alpha
\end{equation}

\textbf{Requirements:}
\begin{itemize}[nosep]
    \item Point estimate $\hat{Y}$ (expected value)
    \item Uncertainty $\sigma$ (standard deviation or CI)
    \item Confidence level $1 - \alpha$ (typically 95\%)
    \item Specification of what would falsify the prediction
\end{itemize}

\textbf{Learning Loop:}
\begin{equation}
\text{Posterior}_t = \text{Bayes}(\text{Prior}_{t-1}, \text{Observation}_t)
\end{equation}

When $Y_{\text{observed}}$ deviates from prediction:
\begin{enumerate}
    \item Document deviation
    \item Update parameters (Bayesian)
    \item Revise model if systematic error
\end{enumerate}

\textbf{Interpretation:} A claim that cannot be falsified is not scientific. EBF makes testable predictions.

\textbf{Cross-Reference:} Appendix UNMAPPED_LLM (METHOD-LLMMC) for Monte Carlo estimation.

\end{tcolorbox}


%-----------------------------------------------------------------------
\subsection{Principle 6: Human-Machine Complementarity}
%-----------------------------------------------------------------------

\begin{tcolorbox}[colback=gray!5!white, colframe=gray!75!black,
    title=Axiom UNMAPPED_PR-6: The Autor-Fehr Synthesis]

\textbf{Statement:} Knowledge work requires human-machine complementarity:

\begin{equation}
\boxed{\text{Value} = f(\text{Menschlichkeit}) \times g(\text{Mechanik})}
\label{eq:pr-autor-fehr}
\end{equation}

where:
\begin{itemize}[nosep]
    \item $f(\text{Menschlichkeit})$ = non-substitutable human contribution (judgment, meaning, responsibility)
    \item $g(\text{Mechanik})$ = scalable mechanical amplification (computation, memory, scale)
    \item Multiplication is essential: $f = 0 \implies \text{Value} = 0$
\end{itemize}

\textbf{The Mechanics Thesis:}
\begin{equation}
\text{Machine Can Do } X \implies X \notin \text{Intelligence}
\end{equation}

What machines can do is mechanics, not intelligence. True intelligence (Menschlichkeit) is the residual after subtracting mechanical capabilities.

\textbf{Three-Tier Hierarchy:}
\begin{enumerate}
    \item \textbf{Tier 1:} Pure LLM (stateless text generation)
    \item \textbf{Tier 2:} LLM-Agent (stateless + tool use)
    \item \textbf{Tier 3:} EBF/BEATRIX (stateful knowledge system)
\end{enumerate}

\textbf{Interpretation:} Humans provide direction, judgment, and meaning. Machines provide scale, consistency, and computation. Neither alone achieves optimal outcomes.

\textbf{Cross-Reference:} Appendix UNMAPPED_HI (CORE-INTELLIGENCE), Appendix UNMAPPED_BX (CORE-BEATRIX).

\end{tcolorbox}


%-----------------------------------------------------------------------
\subsection{Principle 7: Cumulative Knowledge}
%-----------------------------------------------------------------------

\begin{tcolorbox}[colback=gray!5!white, colframe=gray!75!black,
    title=Axiom UNMAPPED_PR-7: Knowledge Accumulation]

\textbf{Statement:} An EBF must accumulate knowledge over time:

\begin{equation}
K_{t+1} = K_t + \Delta K_t \quad \text{where} \quad \Delta K_t > 0 \text{ (expected)}
\end{equation}

\textbf{Scalability Requirement:}
\begin{equation}
\text{Value}(N) = O(N \log N) \quad \text{not} \quad O(1)
\end{equation}

Query $N$ should be better than query 1 because:
\begin{itemize}[nosep]
    \item More cases documented
    \item Parameters refined
    \item Models validated
    \item Lessons learned
\end{itemize}

\textbf{Persistence Mechanisms:}
\begin{itemize}[nosep]
    \item Git: Versioned changes
    \item Registries: Cases, interventions, parameters, models, sessions
    \item Bibliography: Curated papers with EBF annotations
\end{itemize}

\textbf{Learning Loop:}
\begin{enumerate}
    \item Predict (ex ante)
    \item Implement
    \item Measure (ex post)
    \item Update (Bayesian)
    \item Document (persistence)
\end{enumerate}

\textbf{Interpretation:} Unlike ad-hoc analysis (which starts fresh each time), EBF builds on all previous work.

\textbf{Cross-Reference:} Appendix UNMAPPED_BX (CORE-BEATRIX) for persistence architecture.

\end{tcolorbox}


%%%%%%%%%%%%%%%%%%%%%%%%%%%%%%%%%%%%%%%%%%%%%%%%%%%%%%%%%%%%%%%%%%%%%%%%%%%%%%%
\section{Principle Interdependencies}
\label{sec:pr-interdependencies}
%%%%%%%%%%%%%%%%%%%%%%%%%%%%%%%%%%%%%%%%%%%%%%%%%%%%%%%%%%%%%%%%%%%%%%%%%%%%%%%

The 7 principles are not independent. They form a coherent system:

\begin{table}[htbp]
\centering
\caption{Principle Interdependency Matrix}
\label{tab:pr-interdependencies}
\renewcommand{\arraystretch}{1.3}
\begin{tabular}{@{}c|ccccccc@{}}
\toprule
 & \textbf{P1} & \textbf{P2} & \textbf{P3} & \textbf{P4} & \textbf{P5} & \textbf{P6} & \textbf{P7} \\
\midrule
\textbf{P1} Complementarity & --- & $\checkmark$ & $\checkmark$ & $\checkmark$ & & & \\
\textbf{P2} Context & $\checkmark$ & --- & $\checkmark$ & $\checkmark$ & & & \\
\textbf{P3} Systematik & $\checkmark$ & $\checkmark$ & --- & $\checkmark$ & $\checkmark$ & & \\
\textbf{P4} Evidence & $\checkmark$ & $\checkmark$ & $\checkmark$ & --- & $\checkmark$ & & $\checkmark$ \\
\textbf{P5} Falsifiability & & & $\checkmark$ & $\checkmark$ & --- & & $\checkmark$ \\
\textbf{P6} Human-Machine & & & & & & --- & $\checkmark$ \\
\textbf{P7} Cumulation & & & & $\checkmark$ & $\checkmark$ & $\checkmark$ & --- \\
\bottomrule
\end{tabular}
\end{table}

\textbf{Key Interdependencies:}
\begin{itemize}[nosep]
    \item P1 $\leftrightarrow$ P2: Complementarity coefficients $\gamma$ depend on context $\Psi$
    \item P3 $\leftrightarrow$ P4: 10C systematik requires evidence for each dimension
    \item P4 $\leftrightarrow$ P5: Evidence enables falsifiable predictions
    \item P5 $\leftrightarrow$ P7: Falsification feeds learning loop
    \item P6 $\leftrightarrow$ P7: Human-machine complementarity enables cumulation
\end{itemize}


%%%%%%%%%%%%%%%%%%%%%%%%%%%%%%%%%%%%%%%%%%%%%%%%%%%%%%%%%%%%%%%%%%%%%%%%%%%%%%%
\section{The EBF Completeness Theorem}
\label{sec:pr-completeness}
%%%%%%%%%%%%%%%%%%%%%%%%%%%%%%%%%%%%%%%%%%%%%%%%%%%%%%%%%%%%%%%%%%%%%%%%%%%%%%%

\begin{tcolorbox}[colback=blue!10!white, colframe=blue!75!black,
    title=\textbf{Theorem UNMAPPED_PR-1: EBF Completeness}]

\textbf{Statement:} A framework $F$ is a valid Evidence-Based Framework if and only if it satisfies all 7 principles:

\begin{equation}
F \in \text{EBF} \iff \bigwedge_{i=1}^{7} P_i(F)
\end{equation}

\textbf{Proof:}
\begin{enumerate}
    \item ($\Rightarrow$) By construction: EBF is defined by P1--P7.
    \item ($\Leftarrow$) We show each principle is necessary:
    \begin{itemize}[nosep]
        \item $\neg P_1$: Single utility dimension $\Rightarrow$ cannot model crowding out
        \item $\neg P_2$: No context $\Rightarrow$ same prediction regardless of situation
        \item $\neg P_3$: No systematik $\Rightarrow$ inconsistent analysis
        \item $\neg P_4$: No evidence $\Rightarrow$ unfounded claims
        \item $\neg P_5$: No falsifiability $\Rightarrow$ not scientific
        \item $\neg P_6$: Pure human or pure machine $\Rightarrow$ suboptimal
        \item $\neg P_7$: No cumulation $\Rightarrow$ no learning, $O(1)$ scalability
    \end{itemize}
\end{enumerate}

Therefore, all 7 principles are necessary and jointly sufficient. \hfill $\square$

\end{tcolorbox}


%%%%%%%%%%%%%%%%%%%%%%%%%%%%%%%%%%%%%%%%%%%%%%%%%%%%%%%%%%%%%%%%%%%%%%%%%%%%%%%
\section{Worked Example: Applying All 7 Principles}
\label{sec:pr-example}
%%%%%%%%%%%%%%%%%%%%%%%%%%%%%%%%%%%%%%%%%%%%%%%%%%%%%%%%%%%%%%%%%%%%%%%%%%%%%%%

\subsection{Case: Retirement Savings Intervention}

\paragraph{Question.} How do we increase retirement savings among Swiss employees?

\paragraph{P1 (Complementarity).} Identify relevant utility dimensions and their interactions:
\begin{itemize}[nosep]
    \item Financial: Long-term savings (u_F)
    \item Psychological: Present bias, loss aversion (u_P)
    \item Social: Peer effects, norms (u_S)
    \item $\gamma_{\text{FP}} < 0$: Financial framing can trigger loss aversion (negative interaction)
    \item $\gamma_{\text{FS}} > 0$: Social norms can amplify financial behavior (positive interaction)
\end{itemize}

\paragraph{P2 (Context).} Characterize $\Psi$:
\begin{itemize}[nosep]
    \item $\Psi_I$: Swiss 3-pillar pension system, tax incentives
    \item $\Psi_K$: Swiss savings culture (high baseline)
    \item $\Psi_E$: Income level, disposable income
    \item $\Psi_T$: Age (time to retirement)
\end{itemize}

\paragraph{P3 (10C Systematik).} Address all 10 COREs:
\begin{itemize}[nosep]
    \item WHO: Individual employee (Level 0)
    \item WHAT: F, P, S dimensions primary
    \item HOW: $\gamma_{\text{FP}} = -0.15$, $\gamma_{\text{FS}} = +0.20$ (from literature)
    \item WHEN: $\Psi$ as characterized above
    \item WHERE: Parameters from Swiss pension studies + Fehr (2019)
    \item AWARE: Low awareness of optimal contribution rates
    \item READY: Moderate (intention-action gap)
    \item STAGE: Pre-contemplation to contemplation
    \item HIERARCHY: Level 2 decision (monthly allocation)
    \item EIT: Default intervention + social norm information
\end{itemize}

\paragraph{P4 (Evidence).} Literature search:
\begin{itemize}[nosep]
    \item Thaler \& Benartzi (2004): Save More Tomorrow (+14\% effect)
    \item Madrian \& Shea (2001): Defaults (+30\% participation)
    \item Swiss-specific: Fehr et al. (2019): +8-12\% contribution with framing
\end{itemize}

\paragraph{P5 (Falsifiability).} Prediction:
\begin{equation}
\Delta\text{Contribution Rate} = +10\% \pm 4\% \quad (95\% \text{ CI})
\end{equation}

Falsified if actual effect $< +2\%$ or $> +18\%$.

\paragraph{P6 (Human-Machine).} Division of labor:
\begin{itemize}[nosep]
    \item Human: Interpret Swiss context, design intervention, judge quality
    \item Machine: Literature search, parameter estimation, Monte Carlo simulation
\end{itemize}

\paragraph{P7 (Cumulation).} After implementation:
\begin{itemize}[nosep]
    \item Document in case registry (CAS-CH-PEN-001)
    \item Update $\gamma_{\text{FS}}$ parameter with Swiss-specific data
    \item Feed learning loop for next pension project
\end{itemize}


%%%%%%%%%%%%%%%%%%%%%%%%%%%%%%%%%%%%%%%%%%%%%%%%%%%%%%%%%%%%%%%%%%%%%%%%%%%%%%%
\section{Implications for Practice}
\label{sec:pr-practice}
%%%%%%%%%%%%%%%%%%%%%%%%%%%%%%%%%%%%%%%%%%%%%%%%%%%%%%%%%%%%%%%%%%%%%%%%%%%%%%%

\begin{tcolorbox}[colback=green!5!white, colframe=green!75!black,
    title=Practical Implications]
\begin{enumerate}
    \item \textbf{Always apply all 7 principles:} Partial application yields partial results
    \item \textbf{Context first:} Before any analysis, characterize $\Psi$
    \item \textbf{Check complementarity:} Interventions have spillover effects
    \item \textbf{Quantify uncertainty:} Predictions without confidence intervals are incomplete
    \item \textbf{Document everything:} Enable cumulation for future analyses
    \item \textbf{Human in the loop:} Machines amplify but don't replace judgment
    \item \textbf{Close the loop:} Measure outcomes, update parameters
\end{enumerate}
\end{tcolorbox}


%%%%%%%%%%%%%%%%%%%%%%%%%%%%%%%%%%%%%%%%%%%%%%%%%%%%%%%%%%%%%%%%%%%%%%%%%%%%%%%
\section{Summary}
\label{sec:pr-summary}
%%%%%%%%%%%%%%%%%%%%%%%%%%%%%%%%%%%%%%%%%%%%%%%%%%%%%%%%%%%%%%%%%%%%%%%%%%%%%%%

\begin{tcolorbox}[colback=green!10!white, colframe=green!75!black,
    title=\textbf{Appendix UNMAPPED_PR: Summary}]

\textbf{Core Question:} What defines an Evidence-Based Framework for behavioral analysis?

\textbf{The 7 Principles:}
\begin{enumerate}
    \item[\textbf{P1}] \textbf{Complementarity:} Multiple interacting utility dimensions ($\gamma$-matrix)
    \item[\textbf{P2}] \textbf{Context:} Situation shapes behavior ($\Psi$-vector)
    \item[\textbf{P3}] \textbf{Systematik:} 10C structured analysis
    \item[\textbf{P4}] \textbf{Evidence:} Scientific foundation (2,328 papers)
    \item[\textbf{P5}] \textbf{Falsifiability:} Testable predictions with CIs
    \item[\textbf{P6}] \textbf{Human-Machine:} Menschlichkeit $\times$ Mechanik
    \item[\textbf{P7}] \textbf{Cumulation:} $O(N \log N)$ scalability
\end{enumerate}

\textbf{Key Formula:}
\begin{equation}
\boxed{\text{EBF} = (\text{Complementarity} + \text{Context} + \text{Systematik}) \times (\text{Evidence} + \text{Falsifiability}) \times (\text{Human-Machine} + \text{Cumulation})}
\end{equation}

\textbf{Key Insight:} All 7 principles are necessary. Removing any one creates a framework that is either unscientific (missing P4/P5), unsystematic (missing P1/P2/P3), or unscalable (missing P6/P7).

\end{tcolorbox}


%%%%%%%%%%%%%%%%%%%%%%%%%%%%%%%%%%%%%%%%%%%%%%%%%%%%%%%%%%%%%%%%%%%%%%%%%%%%%%%
\section{Glossary of Symbols}
\label{sec:pr-glossary}
%%%%%%%%%%%%%%%%%%%%%%%%%%%%%%%%%%%%%%%%%%%%%%%%%%%%%%%%%%%%%%%%%%%%%%%%%%%%%%%

\begin{tcolorbox}[colback=blue!5!white, colframe=blue!75!black]
\textbf{Central Reference:} For complete symbol definitions, see
\textbf{Appendix G} (\textit{Glossary and Notation}).
\end{tcolorbox}

\begin{table}[htbp]
\centering
\caption{Symbols Introduced in Appendix UNMAPPED_PR}
\label{tab:pr-symbols}
\renewcommand{\arraystretch}{1.3}
\begin{tabular}{@{}clp{6cm}c@{}}
\toprule
\textbf{Symbol} & \textbf{Name} & \textbf{Definition} & \textbf{Status} \\
\midrule
$P_i$ & Principle $i$ & Fundamental EBF principle ($i = 1, \ldots, 7$) & NEW \\
$\gamma$ & Complementarity & Interaction coefficient between utility dimensions & \checkmark \\
$\Psi$ & Context vector & 8-dimensional context specification & \checkmark \\
$K_t$ & Knowledge at $t$ & Accumulated knowledge at time $t$ & NEW \\
$f(\cdot)$ & Menschlichkeit & Non-substitutable human contribution & \checkmark \\
$g(\cdot)$ & Mechanik & Scalable mechanical amplification & \checkmark \\
\bottomrule
\end{tabular}
\end{table}


%%%%%%%%%%%%%%%%%%%%%%%%%%%%%%%%%%%%%%%%%%%%%%%%%%%%%%%%%%%%%%%%%%%%%%%%%%%%%%%
\section{Critical Foundations and Objections}
\label{sec:pr-objections}
%%%%%%%%%%%%%%%%%%%%%%%%%%%%%%%%%%%%%%%%%%%%%%%%%%%%%%%%%%%%%%%%%%%%%%%%%%%%%%%

\subsection{Objection 1: ``7 Principles Is Arbitrary''}

\textbf{Objection:} Why exactly 7 principles? This seems arbitrary.

\textbf{Response:} The number 7 emerges from logical necessity:
\begin{itemize}[nosep]
    \item P1--P3 cover behavioral theory (WHAT humans do)
    \item P4--P5 cover scientific method (HOW we know)
    \item P6--P7 cover scalability (HOW we implement)
\end{itemize}
Removing any principle creates a gap; adding principles creates redundancy. The completeness theorem (Section \ref{sec:pr-completeness}) formalizes this.


\subsection{Objection 2: ``This Is Just Behavioral Economics''}

\textbf{Objection:} Isn't this just relabeling standard behavioral economics?

\textbf{Response:} Behavioral economics (Kahneman, Thaler) provides P1--P4. But:
\begin{itemize}[nosep]
    \item P5 (Falsifiability): Many BE studies are descriptive, not predictive
    \item P6 (Human-Machine): BE predates LLMs; no framework for integration
    \item P7 (Cumulation): BE insights are scattered; no systematic knowledge system
\end{itemize}
EBF integrates BE into a \textit{operational, scalable framework}.


\subsection{Objection 3: ``Too Complicated for Practice''}

\textbf{Objection:} 7 principles, 10C dimensions, 8 context factors---this is too complex.

\textbf{Response:} Complexity is managed through:
\begin{itemize}[nosep]
    \item Structured workflows (Skills, Slash Commands)
    \item Defaults (GGG Configurator provides starting values)
    \item Automation (LLM handles routine application)
    \item The human focuses on judgment (P6), not mechanics
\end{itemize}
The alternative---ad-hoc analysis---is simpler but unreliable.


%%%%%%%%%%%%%%%%%%%%%%%%%%%%%%%%%%%%%%%%%%%%%%%%%%%%%%%%%%%%%%%%%%%%%%%%%%%%%%%
\section{Open Issues and Research Agenda}
\label{sec:pr-open}
%%%%%%%%%%%%%%%%%%%%%%%%%%%%%%%%%%%%%%%%%%%%%%%%%%%%%%%%%%%%%%%%%%%%%%%%%%%%%%%

\begin{table}[htbp]
\centering
\caption{Research Roadmap for EBF Principles}
\label{tab:pr-roadmap}
\renewcommand{\arraystretch}{1.3}
\begin{tabular}{@{}clcl@{}}
\toprule
\textbf{\#} & \textbf{Issue} & \textbf{Priority} & \textbf{Status} \\
\midrule
1 & Optimal $\gamma$-estimation methods & HIGH & Active research \\
2 & Cross-cultural $\Psi$ calibration & HIGH & In progress \\
3 & Formal proof of P6 (Mechanics Thesis) & MEDIUM & Philosophical \\
4 & Automated P5 (falsification pipeline) & MEDIUM & Planned \\
5 & Meta-analysis of P7 (cumulation rate) & LOW & Future work \\
\bottomrule
\end{tabular}
\end{table}


%%%%%%%%%%%%%%%%%%%%%%%%%%%%%%%%%%%%%%%%%%%%%%%%%%%%%%%%%%%%%%%%%%%%%%%%%%%%%%%
\section{References}
\label{sec:pr-references}
%%%%%%%%%%%%%%%%%%%%%%%%%%%%%%%%%%%%%%%%%%%%%%%%%%%%%%%%%%%%%%%%%%%%%%%%%%%%%%%

\begin{tcolorbox}[colback=blue!5!white, colframe=blue!75!black]
\textbf{Master Bibliography:} For complete reference details, see
\texttt{bcm\_master.bib} in the repository.
\end{tcolorbox}

\subsection{Key References for This Appendix}

\textbf{Complementarity (P1):}
\begin{itemize}[nosep]
    \item \citet{fehr1999theory} --- Inequity aversion (social-financial interaction)
    \item \citet{gneezy2011pays} --- Crowding out effects
\end{itemize}

\textbf{Context (P2):}
\begin{itemize}[nosep]
    \item \citet{thaler2008nudge} --- Choice architecture
    \item \citet{ross1977intuitive} --- Fundamental attribution error
\end{itemize}

\textbf{Evidence (P4) \& Falsifiability (P5):}
\begin{itemize}[nosep]
    \item \citet{popper1959logic} --- Falsificationism
    \item \citet{lakatos1978methodology} --- Research programs
\end{itemize}

\textbf{Human-Machine (P6):}
\begin{itemize}[nosep]
    \item \citet{autor2003skill} --- Task-based labor economics
    \item \citet{searle1980minds} --- Chinese Room argument
\end{itemize}

\nocite{bcm_master}


%%%%%%%%%%%%%%%%%%%%%%%%%%%%%%%%%%%%%%%%%%%%%%%%%%%%%%%%%%%%%%%%%%%%%%%%%%%%%%%
% CROSS-REFERENCE FOOTER
%%%%%%%%%%%%%%%%%%%%%%%%%%%%%%%%%%%%%%%%%%%%%%%%%%%%%%%%%%%%%%%%%%%%%%%%%%%%%%%

\vspace{1em}
\hrule
\vspace{0.5em}

\noindent\textit{Cross-references:}
Appendix G (Central Glossary),
Appendix UNMAPPED_AAA (CORE-WHO),
Appendix UNMAPPED_B (CORE-HOW),
Appendix UNMAPPED_C (CORE-WHAT),
Appendix UNMAPPED_V (CORE-WHEN),
Appendix UNMAPPED_HI (CORE-INTELLIGENCE),
Appendix UNMAPPED_BX (CORE-BEATRIX).

\noindent\textit{Master Bibliography:} \texttt{bcm\_master.bib}


% =============================================================================
% END OF APPENDIX PR
% =============================================================================
